\documentclass[conference]{IEEEtran}

\usepackage[T1]{fontenc}
\usepackage{cite}
\usepackage{booktabs}
\usepackage{graphicx}
\usepackage{amsmath}
\usepackage{url}
\usepackage{xcolor}

\title{Beyond Same-Corpus Accuracy: A Cross-Source Study of Hindi Audio Deepfake Detection}

\author{
\IEEEauthorblockN{Prathmesh Pravin Sanap, Shivam Vinod Shewale, Yash Vinayak Jadhav, Snehal Sandeep Sanap}
\IEEEauthorblockA{Department of Computer Engineering\\
K. K. Wagh Institute of Engineering Education and Research, Nashik, India\\
Under the guidance of Prof. S. S. Bhandare}
}

\begin{document}
\maketitle

\begin{abstract}
Audio deepfake detectors can score almost perfectly when their training and
test data come from the same dataset. However, they may fail when the audio
comes from a new source. We study this problem for Hindi speech. We use real
speech from Kathbath and Common Voice, and synthetic speech from IndicSynth
and SEA-Spoof. Early source-held-out experiments produced equal error rates
(EERs) between 3.74\% and 40.46\%. They also showed that the model was
learning differences between datasets, including silence and recording
conditions. We therefore trained the model with real and fake clips that had
the same transcript. This prevents the words in a clip from becoming a
shortcut for the label. Three EfficientNet-B0 models were trained on 3,000
checked SEA-Spoof transcript pairs. The models were almost perfect on
held-out SEA-Spoof data. On unseen sources, mean balanced accuracy was
89.07\% (SD 1.15), ROC-AUC was 0.9501 (SD 0.0205), and EER was 10.98\%
(SD 3.69). Most errors came from Kathbath: only 39.44\% of its real clips
were classified correctly. Accuracy on every other external source was above
98.6\%. The results show why deepfake detectors must be tested on sources
that are absent from training. They also show that matching transcripts does
not remove every dataset-specific effect.
\end{abstract}

\begin{IEEEkeywords}
audio deepfake detection, Hindi speech, cross-source evaluation,
transcript pairing, EfficientNet, source bias
\end{IEEEkeywords}

\section{Introduction}
Modern text-to-speech and voice-conversion systems can create convincing
synthetic speech. This creates risks for voice authentication, media
verification, and digital forensics. A detector must therefore do more than
separate the examples in one benchmark. It should also work when speakers,
recording devices, codecs, and synthesis systems change.

This requirement is difficult to measure. Real and fake clips often come
from different datasets. A model can then learn which dataset produced a
clip instead of learning general signs of synthesis. Random train--test
splits make this problem worse because every source appears in both parts of
the experiment. Prior work has reported large losses when audio deepfake
detectors are moved to unseen data~\cite{muller2022generalize,muller2024harder}.

We investigate this problem for Hindi. Our first pooled experiment achieved
perfect test scores, but source-held-out results varied widely. Signal and
embedding checks suggested that recording and silence patterns were acting
as shortcuts. We then built a larger experiment using transcript-paired
SEA-Spoof data~\cite{wu2025seaspoof}. Each selected transcript had one real
and one fake clip, and all clips with the same transcript were kept in one
partition.

This paper makes three contributions:
\begin{itemize}
    \item a leakage-checked protocol for transcript-paired Hindi deepfake
    detection;
    \item a three-seed cross-source evaluation on genuine and synthetic
    sources that were absent from training; and
    \item a source-level error analysis that identifies Kathbath, rather
    than the unseen fake generator, as the main failure case.
\end{itemize}

We do not claim that transcript pairing is better than ordinary sampling.
That claim needs a count-matched non-paired control, which is left for future
work. Our present claim is narrower: same-corpus scores greatly overstate
performance, and the remaining errors are strongly source dependent.

\section{Related Work}
Audio deepfake detection is often studied using spectral representations or
learned waveform features. Müller et al.~\cite{muller2022generalize} compared
several detection systems under a common setup and showed that performance
can fall sharply outside established benchmarks. Their later analysis found
that cross-domain failure is driven mainly by differences between domains,
not simply by one domain containing harder fakes~\cite{muller2024harder}.
Kawa et al.~\cite{kawa2022attack} also studied attack-independent evaluation
and used disjoint attacks to measure generalization.

Self-supervised speech representations are another strong direction. Pascu
et al.~\cite{pascu2024calibrated} showed that frozen pretrained
representations with a small classifier can generalize well across several
deepfake datasets. Our study instead uses a compact image-based baseline so
that the effect of dataset construction and source shift can be examined
with a simple, reproducible pipeline.

Hindi and other Indian languages have fewer deepfake benchmarks than
English. SEA-Spoof provides paired real and spoof speech for several South
and Southeast Asian languages, including Hindi~\cite{wu2025seaspoof}.
Kathbath is part of IndicSUPERB and contains labelled speech from 12 Indian
languages~\cite{javed2022indicsuperb}. IndicSynth provides synthetic speech
for low-resource Indian languages~\cite{sharma2025indicsynth}. We use these
resources to study transfer between independent Hindi sources.

\section{Data}
Table~\ref{tab:data} summarizes the data used in the final experiment.
SEA-Spoof is the only training source. Kathbath, Common Voice, and IndicSynth
are used only for external evaluation. Common Voice is a crowdsourced,
multilingual speech corpus~\cite{ardila2020commonvoice}.

\begin{table}[t]
\caption{Datasets used in the final experiment}
\label{tab:data}
\centering
\small
\begin{tabular}{llll}
\toprule
Dataset & Label & Role & Samples/seed \\
\midrule
SEA-Spoof Hindi & real+fake & train/val/test & 6,000 total \\
Kathbath Hindi & real & external test & 100 \\
Common Voice male & real & external test & 100 \\
Common Voice female & real & external test & 100 \\
IndicSynth Hindi & fake & external test & 100 \\
\bottomrule
\end{tabular}
\end{table}

\subsection{Transcript-Paired Training Set}
We selected 3,000 Hindi transcripts from the SEA-Spoof training split. For
each transcript, we kept one bona fide clip and one spoof clip. The final
corpus therefore contained 6,000 clips. A protocol audit checked the
following conditions:
\begin{itemize}
    \item all 3,000 transcript groups contained both labels;
    \item no transcript appeared in more than one partition;
    \item no transcript or feature path was missing; and
    \item no file path was duplicated.
\end{itemize}

For each run, transcript groups were shuffled using the training seed and
split 80/10/10. This gave 4,800 training clips, 600 validation clips, and
600 same-corpus test clips. Both labels were exactly balanced in every
partition.

\subsection{External Evaluation Data}
We evaluated each trained model on three deterministic sample manifests.
Each manifest contained 100 clips from each external source. A manifest thus
contained 300 real clips and 100 fake clips. Because the classes were not
balanced, we emphasize balanced accuracy, EER, and ROC-AUC rather than raw
accuracy. The sample seeds were 42, 1337, and 2026. These seeds change the
selected external clips; they do not retrain the model.

\section{Method}
\subsection{Audio Preprocessing}
Every clip was converted to mono at 16~kHz. Silence was removed with a
30~dB threshold, the waveform was peak-normalized, and the result was padded
or truncated to four seconds. We computed a log-mel spectrogram with 128 mel
bands, a 1,024-sample FFT and window, a 512-sample hop, and a frequency range
of 0--8~kHz. Each feature map was resized to $224\times224$.

\subsection{Classifier and Training}
The classifier was an ImageNet-pretrained EfficientNet-B0~\cite{tan2019efficientnet}
adapted for a single-channel spectrogram and two output classes. We did not
use the optional BiLSTM in the reported experiments. The backbone was
fine-tuned with AdamW, a learning rate of $3\times10^{-4}$, weight decay of
$10^{-5}$, batch size 32, cosine learning-rate scheduling, mixed precision,
and gradient clipping at 5.0. Training was limited to 50 epochs and stopped
after seven epochs without improvement in validation F1. The best
validation-F1 checkpoint was used for evaluation.

We trained three models with seeds 42, 1337, and 2026. The seed controlled
the transcript split, initial weights, and batch order. This is more useful
than reporting a single favorable run because it measures sensitivity to
random training choices.

\subsection{Metrics}
We report accuracy, balanced accuracy, precision, recall, F1, Matthews
correlation coefficient (MCC), ROC-AUC, and EER. EER is the operating point
at which the false-positive and false-negative rates are equal. For each
model, we first averaged the three external sample runs. We then calculated
the mean and sample standard deviation across the three independently
trained models.

\section{Pilot Source-Shift Analysis}
Before the final experiment, we trained four models using different pairs of
real and fake sources. Each model was tested on the complementary sources.
Table~\ref{tab:factorial} shows that performance depended strongly on which
sources were used for training. These were single-run pilot experiments and
are included as motivation, not as the final comparison.

\begin{table}[t]
\caption{Pilot source-held-out experiments}
\label{tab:factorial}
\centering
\small
\begin{tabular}{llrr}
\toprule
Training real & Training fake & ROC-AUC & EER \\
\midrule
Kathbath & IndicSynth & 0.9939 & 0.0374 \\
Common Voice & IndicSynth & 0.8165 & 0.2325 \\
Common Voice & SEA-Spoof & 0.8761 & 0.2160 \\
Kathbath & SEA-Spoof & 0.6273 & 0.4046 \\
\bottomrule
\end{tabular}
\end{table}

The worst model had almost equal embedding separation for source identity
and class label. Signal checks also found large source differences in
silence and trailing padding. These checks do not prove causation, but they
explain why a source-resolved final evaluation was necessary.

\section{Results}
\subsection{Same-Corpus Evaluation}
All three models were almost perfect on held-out SEA-Spoof data. Test
accuracy ranged from 0.9950 to 1.0000. This confirms that the paired task is
learnable without transcript leakage. It does not show that the detector
will generalize to another dataset.

\subsection{Cross-Source Evaluation}
Table~\ref{tab:main-results} reports the result for each trained model. Each
row is the average of that model's three external manifests. The final row
reports variation across model-training seeds.

\begin{table*}[t]
\caption{Cross-source results. Mean $\pm$ SD is calculated across three independently trained models.}
\label{tab:main-results}
\centering
\begin{tabular}{lrrrrrr}
\toprule
Training seed & Accuracy & Balanced acc. & F1 & ROC-AUC & EER & MCC \\
\midrule
42   & 0.8625 & 0.9039 & 0.7825 & 0.9625 & 0.0922 & 0.7210 \\
1337 & 0.8242 & 0.8828 & 0.7400 & 0.9615 & 0.0850 & 0.6706 \\
2026 & 0.8400 & 0.8856 & 0.7534 & 0.9264 & 0.1522 & 0.6827 \\
\midrule
Mean $\pm$ SD & $0.8422\pm0.0193$ & $0.8907\pm0.0115$ & $0.7586\pm0.0217$ & $0.9501\pm0.0205$ & $0.1098\pm0.0369$ & $0.6914\pm0.0264$ \\
\bottomrule
\end{tabular}
\end{table*}

The same-corpus EER was near zero, while cross-source EER increased to
10.98\% on average. Model seed 2026 was weaker than the other models, with
15.22\% EER. This shows why results from one seed can be misleading.

\subsection{Source-Level Errors}
Table~\ref{tab:source-results} shows that the loss was not spread evenly
across datasets. The models were almost perfect on both Common Voice subsets
and on unseen IndicSynth fake speech. In contrast, mean accuracy on Kathbath
was only 39.44\%.

\begin{table}[t]
\caption{Accuracy by external source}
\label{tab:source-results}
\centering
\small
\begin{tabular}{llrr}
\toprule
Source & Label & Mean acc. & Manifest range \\
\midrule
Kathbath & real & 0.3944 & 0.28--0.53 \\
Common Voice male & real & 1.0000 & 1.00--1.00 \\
Common Voice female & real & 0.9867 & 0.97--0.99 \\
IndicSynth & fake & 0.9878 & 0.95--1.00 \\
\bottomrule
\end{tabular}
\end{table}

The aggregate accuracy is therefore mainly reduced by real Kathbath clips
being labelled as fake. The unseen fake generator was not the main problem.
This result is consistent with the earlier observation that recording and
silence properties can become shortcuts. The high ROC-AUC alongside lower
fixed-threshold F1 also suggests a calibration shift: score ranking remains
useful even when a fixed decision threshold produces many errors.

\section{Discussion}
The experiment gives two clear results. First, same-corpus performance is an
optimistic upper bound. A model can be perfect on held-out examples from the
same source and still fail on a new genuine-speech corpus. Second,
cross-source failure must be reported separately for each source. An
aggregate score alone would hide the fact that almost all errors came from
Kathbath.

Transcript pairing removes one possible shortcut: the spoken text cannot be
associated with only one class. It does not remove microphones, codecs,
speakers, silence patterns, or generator artifacts. The results therefore
show that content control is useful for building a cleaner experiment, but
it is not a complete solution to domain shift.

The present study does not measure the causal improvement from pairing. The
classifier receives individual spectrograms and does not use a pairwise loss.
A fair comparison requires an equally sized training set whose real and fake
transcripts are not matched, while keeping the architecture, optimization,
and external manifests fixed. We leave this control to future work rather
than claiming an unsupported improvement.

\section{Limitations}
This study has several limitations. It covers Hindi only. Speaker IDs were
not preserved in the metadata, so speaker-disjoint splits could not be
verified. The external manifests contain 100 clips per source and overlap
across sampling seeds; a fixed exhaustive test set would be stronger. We
used three training seeds and one model family. No count-matched non-paired
control or normalization ablation was completed. We also did not compute
minimum tandem detection cost because the experiment did not include a
compatible automatic speaker verification system. Finally, Kathbath
performance is too low for the detector to be used as a reliable real-world
decision system.

\section{Ethical Considerations}
Audio deepfake detection can support media verification and fraud analysis,
but its output is not proof that a person acted deceptively. False positives
can wrongly label genuine speech as fake, as the Kathbath results clearly
show. Such a detector should not be the sole decision maker in a legal,
financial, or identity-verification setting. Dataset licenses and consent
conditions must be followed, and restricted audio should not be
redistributed with the code.

\section{Conclusion}
Three transcript-paired Hindi deepfake detectors were almost perfect on
held-out SEA-Spoof data but reached 10.98\% $\pm$ 3.69\% EER on external
sources. ROC-AUC remained 0.9501 $\pm$ 0.0205. The loss was concentrated on
Kathbath genuine speech, while Common Voice and IndicSynth accuracy stayed
above 98.6\%. These findings show that same-corpus accuracy can hide a large,
source-specific weakness. Future work should add a count-matched non-paired
control, test silence and loudness normalization, retain speaker IDs, and
compare EfficientNet-B0 with a speech-specific baseline.

\bibliographystyle{IEEEtran}
\bibliography{references}

\end{document}
